\documentclass[11pt,a4paper]{article}

\usepackage[utf8]{inputenc}
\usepackage[T1]{fontenc}
\usepackage[spanish]{babel}
\usepackage{amsmath,amssymb,mathtools}
\usepackage{geometry}
\usepackage{enumitem}
\usepackage{array}
\usepackage{lmodern}
\usepackage{microtype}

\geometry{margin=2.4cm}
\setlength{\parindent}{0pt}
\setlength{\parskip}{0.5em}

\newcommand{\R}{\mathbb{R}}
\newcommand{\K}{\mathbb{K}}
\newcommand{\Ker}{\operatorname{Ker}}
\newcommand{\Ima}{\operatorname{Im}}
\newcommand{\rg}{\operatorname{rg}}
\newcommand{\Span}{\operatorname{L}}
\newcommand{\Mat}{\operatorname{M}}
\newcommand{\Id}{\operatorname{Id}}
\newcommand{\an}{\operatorname{an}}

\begin{document}

\begin{center}
	{\LARGE\bfseries Álgebra Lineal}\\[0.4em]
	{\Large Tema 2: Aplicaciones Lineales --- Soluciones}
\end{center}

\section{Ejercicio 1}
	\subsection{Solución}

	Sean $u=(x,y,z)$, $v=(x',y',z')$ y $\alpha,\beta\in\R$.

	\[
	f(\alpha u+\beta v)
	=
	(\alpha z+\beta z',
	\alpha x+\beta x'+\alpha y+\beta y',
	0,
	-\alpha z-\beta z')
	\]

	\[
	=
	\alpha(z,x+y,0,-z)
	+
	\beta(z',x'+y',0,-z')
	\]

	\[
	=
	\alpha f(u)+\beta f(v)
	\]

	Por tanto:

	\[
	\boxed{
	f\text{ es una aplicación lineal}
	}
	\]

\section{Ejercicio 2}
	\subsection{Solución}

	Una aplicación lineal debe verificar:

	\[
	f(0,0,0)=(0,0)
	\]

	Pero:

	\[
	f(0,0,0)=(a,0)
	\]

	Por tanto:

	\[
	\boxed{
	f\text{ es lineal si y sólo si }a=0
	}
	\]

\section{Ejercicio 3}
	\subsection{Solución}

	Sean $u=(x,y,z)$, $v=(x',y',z')$ y $\alpha,\beta\in\R$.

	\[
	f(\alpha u+\beta v)
	=
	\begin{pmatrix}
	a(\alpha x+\beta x')&
	b(\alpha y+\beta y')\\
	c(\alpha z+\beta z')&
	\alpha(x+y+z)+\beta(x'+y'+z')
	\end{pmatrix}
	\]

	\[
	=
	\alpha f(u)+\beta f(v)
	\]

	para cualquier $a,b,c\in\R$.

	\[
	\boxed{
	f\text{ es lineal para todo }a,b,c\in\R
	}
	\]

\section{Ejercicio 4}
	\subsection{Solución}

	Para el núcleo:

	\[
	f(x_1,x_2,x_3)=(0,0,0)
	\]

	implica:

	\[
	x_1+x_2+x_3=0
	\]

	\[
	x_1+x_2-x_3=0
	\]

	\[
	x_3=0
	\]

	Por tanto:

	\[
	x_2=-x_1
	\]

	\[
	\boxed{
	\Ker f
	=
	\Span\{(1,-1,0)\}
	}
	\]

	Como:

	\[
	\dim(\Ker f)=1
	\]

	se tiene:

	\[
	\dim(\Ima f)=3-1=2
	\]

	Además:

	\[
	f(1,0,0)=(1,1,0)
	\]

	\[
	f(0,0,1)=(1,-1,1)
	\]

	son linealmente independientes.

	\[
	\boxed{
	\Ima f
	=
	\Span\{(1,1,0),(1,-1,1)\}
	}
	\]

	Para:

	\[
	V=
	\{(x_1,x_2,-x_1-x_2)\}
	\]

	se tiene:

	\[
	f(x_1,x_2,-x_1-x_2)
	=
	(0,2x_1+2x_2,-x_1-x_2)
	\]

	\[
	=
	(x_1+x_2)(0,2,-1)
	\]

	Por tanto:

	\[
	\boxed{
	f(V)=\Span\{(0,2,-1)\}
	}
	\]

\section{Ejercicio 5}
	\subsection{Solución}

	Para el núcleo:

	\[
	x_1+2x_2+3x_3=0
	\]

	\[
	\boxed{
	\Ker f
	=
	\{(-2s-3t,s,t):s,t\in\R\}
	}
	\]

	\[
	=
	\Span\{(-2,1,0),(-3,0,1)\}
	\]

	Así:

	\[
	\boxed{
	\dim(\Ker f)=2
	}
	\]

	La imagen está generada por:

	\[
	f(1,0,0)=(1,2)
	\]

	\[
	f(0,1,0)=(2,4)=2(1,2)
	\]

	\[
	f(0,0,1)=(3,6)=3(1,2)
	\]

	Por tanto:

	\[
	\boxed{
	\Ima f=\Span\{(1,2)\}
	}
	\]

	\[
	\boxed{
	\dim(\Ima f)=1
	}
	\]

	y:

	\[
	2+1=3=\dim\R^3
	\]

\section{Ejercicio 6}
	\subsection{Solución}

	Para el núcleo:

	\[
	x_3=0,
	\qquad
	x_1+x_2=0
	\]

	\[
	\boxed{
	\Ker f
	=
	\Span\{(1,-1,0)\}
	}
	\]

	La imagen está generada por:

	\[
	f(1,0,0)=(0,1,0)
	\]

	\[
	f(0,0,1)=(1,0,-1)
	\]

	\[
	\boxed{
	\Ima f
	=
	\Span\{(0,1,0),(1,0,-1)\}
	}
	\]

	Además:

	\[
	\boxed{
	f^{-1}(0,0,0)
	=
	\Ker f
	}
	\]

	Para:

	\[
	f(x_1,x_2,x_3)=(1,2,3)
	\]

	se exigiría simultáneamente:

	\[
	x_3=1
	\]

	y:

	\[
	-x_3=3
	\]

	lo cual es imposible.

	\[
	\boxed{
	f^{-1}(1,2,3)=\emptyset
	}
	\]

	Para:

	\[
	f(x_1,x_2,x_3)=(1,0,-1)
	\]

	se obtiene:

	\[
	x_3=1,
	\qquad
	x_1+x_2=0
	\]

	\[
	\boxed{
	f^{-1}(1,0,-1)
	=
	\{(s,-s,1):s\in\R\}
	}
	\]

\section{Ejercicio 7}
	\subsection{Solución}

	Buscamos:

	\[
	1-x+5x^2
	=
	\alpha(1+x^2)+\beta x+\gamma(3+x)
	\]

	Comparando coeficientes:

	\[
	\alpha+3\gamma=1
	\]

	\[
	\beta+\gamma=-1
	\]

	\[
	\alpha=5
	\]

	Por tanto:

	\[
	\alpha=5,
	\qquad
	\beta=\frac13,
	\qquad
	\gamma=-\frac43
	\]

	Aplicando linealidad:

	\[
	f(1-x+5x^2)
	=
	5f(1+x^2)
	+
	\frac13f(x)
	-
	\frac43f(3+x)
	\]

	\[
	=
	5(3+2x^2)
	+
	\frac13(1+x+2x^2)
	+
	\frac43
	\]

	\[
	\boxed{
	f(1-x+5x^2)
	=
	\frac{50}{3}
	+
	\frac13x
	+
	\frac{32}{3}x^2
	}
	\]

\section{Ejercicio 8}
	\subsection{Solución}

	Definimos:

	\[
	w_1=3u_1-2u_2
	\]

	\[
	w_2=4u_1-3u_2
	\]

	Se obtiene:

	\[
	u_1=3w_1-2w_2
	\]

	\[
	u_2=4w_1-3w_2
	\]

	Entonces:

	\[
	g(u_1)
	=
	3g(w_1)-2g(w_2)
	\]

	\[
	=
	7v_1+8v_2-7v_3
	\]

	y:

	\[
	g(u_2)
	=
	4g(w_1)-3g(w_2)
	\]

	\[
	=
	9v_1+9v_2-9v_3
	\]

	Por tanto:

	\[
	\boxed{
	M_{B,B'}(g)
	=
	\begin{pmatrix}
	7&9\\
	8&9\\
	-7&-9
	\end{pmatrix}
	}
	\]

	El núcleo viene dado por:

	\[
	\begin{pmatrix}
	7&9\\
	8&9\\
	-7&-9
	\end{pmatrix}
	\begin{pmatrix}
	x\\y
	\end{pmatrix}
	=
	0
	\]

	que únicamente tiene la solución:

	\[
	x=y=0
	\]

	Por tanto:

	\[
	\boxed{
	g\text{ es inyectiva}
	}
	\]

	Como:

	\[
	\rg M=2<3
	\]

	\[
	\boxed{
	g\text{ no es suprayectiva}
	}
	\]

\section{Ejercicio 9}
	\subsection{Solución}

	Calculamos:

	\[
	f(1,1,1)=(9,11)
	\]

	\[
	f(1,1,0)=(4,4)
	\]

	\[
	f(0,1,1)=(8,9)
	\]

	Expresamos estas imágenes en $B'$:

	\[
	(9,11)
	=
	17(1,1)-2(4,3)
	\]

	\[
	(4,4)
	=
	4(1,1)
	\]

	\[
	(8,9)
	=
	12(1,1)-(4,3)
	\]

	Por tanto:

	\[
	\boxed{
	M_{A',B'}(f)
	=
	\begin{pmatrix}
	17&4&12\\
	-2&0&-1
	\end{pmatrix}
	}
	\]

\section{Ejercicio 10}
	\subsection{Solución}

	Llamamos:

	\[
	w_1=u_1+3u_2+2u_3
	\]

	\[
	w_2=u_2+u_3
	\]

	\[
	w_3=u_1+u_3
	\]

	La matriz respecto de la base $\{w_1,w_2,w_3\}$ en el dominio es:

	\[
	\begin{pmatrix}
	1&-1&4\\
	3&1&2
	\end{pmatrix}
	\]

	El cambio de base da:

	\[
	M_{A',B'}(f)
	=
	\frac12
	\begin{pmatrix}
	1&-1&4\\
	3&1&2
	\end{pmatrix}
	\begin{pmatrix}
	-1&1&-1\\
	3&-1&3\\
	1&-1&1
	\end{pmatrix}
	\]

	\[
	\boxed{
	M_{A',B'}(f)
	=
	\begin{pmatrix}
	4&-1&0\\
	1&0&1
	\end{pmatrix}
	}
	\]

\section{Ejercicio 11}
	\subsection{Solución}

	La matriz de cambio de $C'$ a $C$ es:

	\[
	P=
	\begin{pmatrix}
	1&1&0\\
	1&0&-1\\
	0&1&2
	\end{pmatrix}
	\]

	y:

	\[
	P^{-1}
	=
	\begin{pmatrix}
	-1&2&1\\
	2&-2&-1\\
	-1&1&1
	\end{pmatrix}
	\]

	Como se trata de un endomorfismo:

	\[
	M_{C'}(f)
	=
	P^{-1}M_C(f)P
	\]

	Por tanto:

	\[
	\boxed{
	M_{C'}(f)
	=
	\begin{pmatrix}
	1&0&0\\
	0&1&0\\
	0&0&2
	\end{pmatrix}
	}
	\]

\section{Ejercicio 12}
	\subsection{Solución}

	Una base de:

	\[
	S=
	\{(x,y,z):x+y+z=0\}
	\]

	es:

	\[
	\{(1,0,-1),(0,1,-1)\}
	\]

	y:

	\[
	f(1,0,-1)=(1,0,-1)
	\]

	\[
	f(0,1,-1)=(0,1,-1)
	\]

	Además:

	\[
	\Ker f=S^\perp
	\]

	Por tanto:

	\[
	\Ker f
	=
	\Span\{(1,1,1)\}
	\]

	y:

	\[
	f(1,1,1)=0
	\]

	Así:

	\[
	\boxed{
	M_{B,B_c}(f)
	=
	\begin{pmatrix}
	1&0&0\\
	0&1&0\\
	-1&-1&0
	\end{pmatrix}
	}
	\]

	La matriz en la base canónica es:

	\[
	\boxed{
	M_{B_c}(f)
	=
	\frac13
	\begin{pmatrix}
	2&-1&-1\\
	-1&2&-1\\
	-1&-1&2
	\end{pmatrix}
	}
	\]

	Por tanto:

	\[
	\boxed{
	f(x,y,z)
	=
	\frac13
	(2x-y-z,\;-x+2y-z,\;-x-y+2z)
	}
	\]

	Sea:

	\[
	W=\Span\{(1,0,0)\}
	\]

	El sistema:

	\[
	f(x,y,z)=(1,0,0)
	\]

	es incompatible.

	Por tanto:

	\[
	\boxed{
	f^{-1}(W)=\emptyset
	}
	\]

	Además:

	\[
	\dim(\Ker f)=1
	\]

	\[
	\dim(\Ima f)=2
	\]

	Así:

	\[
	\boxed{
	f\text{ no es inyectiva ni suprayectiva}
	}
	\]

\section{Ejercicio 13}
	\subsection{Solución}

	Calculamos las imágenes de los vectores de $B$:

	\[
	f(1,3,0)=(3,11)
	\]

	\[
	f(1,0,2)=(-1,9)
	\]

	\[
	f(0,4,-2)=(4,-4)
	\]

	Expresándolas en $B'$:

	\[
	(3,11)
	=
	-\frac{35}{2}(2,1)
	+
	\frac{19}{2}(4,3)
	\]

	\[
	(-1,9)
	=
	-\frac{39}{2}(2,1)
	+
	\frac{19}{2}(4,3)
	\]

	\[
	(4,-4)
	=
	-6(2,1)+4(4,3)
	\]

	Por tanto:

	\[
	\boxed{
	M_{B,B'}(f)
	=
	\begin{pmatrix}
	-\frac{35}{2}&-\frac{39}{2}&-6\\
	\frac{19}{2}&\frac{19}{2}&4
	\end{pmatrix}
	}
	\]

\section{Ejercicio 14}
	\subsection{Solución}

	Como:

	\[
	\Ker f
	=
	\{(x,y,z):x-y+2z=0\}
	\]

	y:

	\[
	h=g\circ f
	\]

	se tiene:

	\[
	\Ker f\subseteq\Ker h
	\]

	Una base del núcleo es:

	\[
	(1,1,0),
	\qquad
	(-2,0,1)
	\]

	Además:

	\[
	h(1,1,1)=(-1,0,2)
	\]

	Expresando los vectores canónicos respecto de esta base ampliada se obtiene:

	\[
	\boxed{
	M_{B_1}(h)
	=
	\begin{pmatrix}
	-\frac12&\frac12&-1\\
	0&0&0\\
	1&-1&2
	\end{pmatrix}
	}
	\]

	Para $W$, resolvemos:

	\[
	W=
	\Span\{(-7,-2,1)\}_{B_1}
	\]

	Entonces:

	\[
	h(-7,-2,1)
	=
	\left(
	\frac32,0,-3
	\right)
	\]

	Por tanto:

	\[
	\boxed{
	h(W)
	=
	\Span
	\left\{
	\left(\frac32,0,-3\right)
	\right\}
	}
	\]

	El cambio a la base $B'$ proporciona:

	\[
	\boxed{
	M_{B'}(h)
	=
	\frac14
	\begin{pmatrix}
	5&-3&-2\\
	-5&3&2\\
	5&-3&-2
	\end{pmatrix}
	}
	\]

	Si las ecuaciones propuestas fueran las de $\Ima g$, entonces:

	\[
	\Ima g
	=
	\Span\{(-11,4,1)\}
	\]

	pero:

	\[
	h(1,1,1)=(-1,0,2)
	\]

	pertenece necesariamente a $\Ima g$ y no verifica esas ecuaciones.

	Por tanto:

	\[
	\boxed{
	\text{No pueden ser ecuaciones implícitas de }\Ima g
	}
	\]

\section{Ejercicio 15}
	\subsection{Solución}

	Escribimos los vectores de $B$ por columnas:

	\[
	A=
	\begin{pmatrix}
	1&1&1\\
	1&1&0\\
	1&0&0
	\end{pmatrix}
	\]

	Entonces:

	\[
	A^{-1}
	=
	\begin{pmatrix}
	0&0&1\\
	0&1&-1\\
	1&-1&0
	\end{pmatrix}
	\]

	Las filas de $A^{-1}$ son las formas de la base dual.

	Por tanto:

	\[
	\boxed{
	B^*
	=
	\{
	u_1^*=x_3,\;
	u_2^*=x_2-x_3,\;
	u_3^*=x_1-x_2
	\}
	}
	\]

\section{Ejercicio 16}
	\subsection{Solución}

	Como:

	\[
	f=3v_1^*+4v_2^*
	\]

	se tiene:

	\[
	f(v_1)=3
	\]

	\[
	f(v_2)=4
	\]

	Escribimos:

	\[
	(x_1,x_2)
	=
	\alpha(1,2)+\beta(3,1)
	\]

	de donde:

	\[
	\alpha
	=
	\frac{3x_2-x_1}{5}
	\]

	\[
	\beta
	=
	\frac{2x_1-x_2}{5}
	\]

	Entonces:

	\[
	f(x_1,x_2)
	=
	3\alpha+4\beta
	\]

	\[
	=
	x_1+x_2
	\]

	\[
	\boxed{
	f(x_1,x_2)=x_1+x_2
	}
	\]

\section{Ejercicio 17}
	\subsection{Solución}

	Sea:

	\[
	p(x)=a+bx+cx^2
	\]

	Entonces:

	\[
	p(1)=a+b+c
	\]

	\[
	p'(1)=b+2c
	\]

	\[
	p''(1)=2c
	\]

	\[
	p'''(1)=0
	\]

	Por tanto:

	\[
	\boxed{
	M_{B_1,B_2}(f_\beta)
	=
	\begin{pmatrix}
	\beta&\beta&\beta\\
	0&1&2\\
	0&0&2\\
	0&0&0
	\end{pmatrix}
	}
	\]

	Si:

	\[
	\beta\neq0
	\]

	el rango es $3$, luego:

	\[
	\boxed{
	\dim(\Ima f_\beta)=3,
	\qquad
	\dim(\Ker f_\beta)=0
	}
	\]

	Si:

	\[
	\beta=0
	\]

	el rango es $2$, luego:

	\[
	\boxed{
	\dim(\Ima f_0)=2,
	\qquad
	\dim(\Ker f_0)=1
	}
	\]

	Para $\beta=0$:

	\[
	\Ker f_0
	=
	\{(x,0,0):x\in\R\}_{B_1}
	\]

	\[
	\boxed{
	\Ker f_0=\Span\{(1,0,0)\}
	}
	\]

	y:

	\[
	\Ima f_0
	=
	\{(x_1,x_2,x_3,x_4)_{B_2}:x_1=0,\ x_4=0\}
	\]

	Para $\beta=1$, una matriz:

	\[
	f_1(a+bx+cx^2)
	=
	\begin{pmatrix}
	a+b+c&b+2c\\
	2c&0
	\end{pmatrix}
	\]

	es simétrica si:

	\[
	b+2c=2c
	\]

	\[
	b=0
	\]

	Por tanto:

	\[
	\boxed{
	f_1^{-1}(S_{2\times2})
	=
	\Span\{1,x^2\}
	}
	\]

	Con las bases $B_1'$ y $B_2'$:

	\[
	\boxed{
	M_{B_1',B_2'}(f_1)
	=
	\begin{pmatrix}
	1&0&2\\
	0&-1&-2\\
	0&0&-2\\
	0&0&0
	\end{pmatrix}
	}
	\]

\section{Ejercicio 18}
	\subsection{Solución}

	Los polinomios con término independiente nulo verifican:

	\[
	f(x)=x
	\]

	\[
	f(x^2)=x^2
	\]

	Además:

	\[
	1+x+x^2\in\Ker f
	\]

	por lo que:

	\[
	f(1)+f(x)+f(x^2)=0
	\]

	\[
	f(1)=-x-x^2
	\]

	Por tanto:

	\[
	\boxed{
	M_{B_c}(f)
	=
	\begin{pmatrix}
	0&0&0\\
	-1&1&0\\
	-1&0&1
	\end{pmatrix}
	}
	\]

	El subespacio $S$ tiene base:

	\[
	\{(1,1,1),(1,-1,0)\}
	\]

	y:

	\[
	f(1,1,1)=0
	\]

	\[
	f(1,-1,0)=(0,-2,-1)
	\]

	Por tanto:

	\[
	\boxed{
	f(S)
	=
	\Span\{(0,-2,-1)\}
	}
	\]

	El rango de la matriz es:

	\[
	\rg M=2
	\]

	Por tanto:

	\[
	\dim(\Ima f)=2
	\]

	\[
	\dim(\Ker f)=1
	\]

	Así:

	\[
	\boxed{
	f\text{ no es inyectiva ni suprayectiva}
	}
	\]

	En la base:

	\[
	B=\{1+x+x^2,1+x,1\}
	\]

	se obtiene:

	\[
	M_B(f)
	=
	\begin{pmatrix}
	0&-1&-1\\
	0&1&0\\
	0&0&1
	\end{pmatrix}
	\]

	Calculando las coordenadas del vector dado en $B$ y aplicando la matriz se obtiene:

	\[
	\boxed{
	f(1-3x+5x^2)
	=
	-4x+4x^2
	}
	\]

\section{Ejercicio 19}
	\subsection{Solución}

	La matriz formada por los vectores de $B$ es:

	\[
	A=
	\begin{pmatrix}
	1&0&1\\
	1&1&0\\
	0&1&1
	\end{pmatrix}
	\]

	\[
	A^{-1}
	=
	\frac12
	\begin{pmatrix}
	1&1&-1\\
	-1&1&1\\
	1&-1&1
	\end{pmatrix}
	\]

	Por tanto:

	\[
	\boxed{
	B^*
	=
	\left\{
	\frac12(x_1+x_2-x_3),
	\frac12(-x_1+x_2+x_3),
	\frac12(x_1-x_2+x_3)
	\right\}
	}
	\]

\section{Ejercicio 20}
	\subsection{Solución}

	Del núcleo:

	\[
	z=-2t
	\]

	\[
	y=-2x
	\]

	Por tanto:

	\[
	\Ker f
	=
	\Span\{(1,-2,0,0),(0,0,-2,1)\}
	\]

	Usando los datos:

	\[
	f(e_1)=(2,0,2,0)
	\]

	\[
	f(e_4)=(2,0,0,0)
	\]

	y expresando $e_2,e_3$ mediante $e_1,e_4$ y los vectores del núcleo:

	\[
	f(e_2)=(1,0,1,0)
	\]

	\[
	f(e_3)=(1,0,0,0)
	\]

	Así:

	\[
	\boxed{
	M_{B_c}(f)
	=
	\begin{pmatrix}
	2&1&1&2\\
	0&0&0&0\\
	2&1&0&0\\
	0&0&0&0
	\end{pmatrix}
	}
	\]

	Una base de:

	\[
	V=\{(x,y,z,t):x+y+z+t=0\}
	\]

	es:

	\[
	(1,0,0,-1),
	\quad
	(0,1,0,-1),
	\quad
	(0,0,1,-1)
	\]

	Sus imágenes son:

	\[
	(0,0,2,0),
	\quad
	(-1,0,1,0),
	\quad
	(-1,0,0,0)
	\]

	y tienen rango $2$.

	Por tanto:

	\[
	\boxed{
	f(V)
	=
	\Span\{(-1,0,-1,0),(-1,0,0,0)\}
	}
	\]

	Para la base $B$ del enunciado:

	\[
	\boxed{
	M_B(f)
	=
	\frac12
	\begin{pmatrix}
	3&1&3&-1\\
	3&1&3&-1\\
	3&1&0&0\\
	3&1&0&0
	\end{pmatrix}
	}
	\]

\section{Ejercicio 21}
	\subsection{Solución}

	Calculamos:

	\[
	h(e_1)
	=
	g(u_1-u_2)
	=
	2c_3
	\]

	\[
	h(e_2)
	=
	g(u_2)
	=
	c_1-c_2
	\]

	\[
	h(e_3)
	=
	2g(u_2)
	=
	2c_1-2c_2
	\]

	Por tanto:

	\[
	\boxed{
	M_{B_A,B_C}(h)
	=
	\begin{pmatrix}
	0&1&2\\
	0&-1&-2\\
	2&0&0
	\end{pmatrix}
	}
	\]

	Para:

	\[
	h(x_1,x_2,x_3)=(1,1,1)
	\]

	se obtiene:

	\[
	x_2+2x_3=1
	\]

	\[
	-x_2-2x_3=1
	\]

	que es incompatible.

	\[
	\boxed{
	h^{-1}(1,1,1)=\emptyset
	}
	\]

	Para el núcleo:

	\[
	x_1=0
	\]

	\[
	x_2=-2x_3
	\]

	\[
	\boxed{
	\Ker h
	=
	\Span\{(0,-2,1)\}
	}
	\]

	Para la intersección:

	\[
	V\cap W
	=
	\Span\{(5,1,-2)\}
	\]

	y:

	\[
	h(5,1,-2)=(-3,3,10)
	\]

	Por tanto:

	\[
	\boxed{
	h(V\cap W)
	=
	\Span\{(-3,3,10)\}
	}
	\]

\section{Ejercicio 22}
	\subsection{Solución}

	Para $f$:

	\[
	f(x)
	=
	\begin{pmatrix}
	1&0\\
	0&1
	\end{pmatrix}
	\]

	\[
	f(1)
	=
	\begin{pmatrix}
	0&1\\
	1&0
	\end{pmatrix}
	\]

	Por tanto:

	\[
	\boxed{
	M_{B,B'}(f)
	=
	\begin{pmatrix}
	1&0\\
	0&1\\
	1&0
	\end{pmatrix}
	}
	\]

	Para $g$:

	\[
	\boxed{
	M_{B',B''}(g)
	=
	\begin{pmatrix}
	1&0&0\\
	0&0&1\\
	1&0&1
	\end{pmatrix}
	}
	\]

	Por composición:

	\[
	\boxed{
	M_{B,B''}(g\circ f)
	=
	\begin{pmatrix}
	1&0\\
	1&0\\
	2&0
	\end{pmatrix}
	}
	\]

	Como:

	\[
	V=\Span\{x+1\}
	\]

	\[
	(g\circ f)(x+1)=(1,1,2)
	\]

	\[
	\boxed{
	(g\circ f)(V)
	=
	\Span\{(1,1,2)\}
	}
	\]

	Para $g$:

	\[
	\Ker g
	=
	\left\{
	\begin{pmatrix}
	0&b\\
	b&0
	\end{pmatrix}
	:b\in\R
	\right\}
	\]

	\[
	\boxed{
	\dim(\Ker g)=1
	}
	\]

	Por rango-nulidad:

	\[
	\boxed{
	\dim(\Ima g)=2
	}
	\]

	Por tanto:

	\[
	\boxed{
	g\text{ no es inyectiva ni suprayectiva}
	}
	\]

	En la base:

	\[
	B^*=
	\left\{
	\begin{pmatrix}1&0\\0&1\end{pmatrix},
	\begin{pmatrix}0&1\\1&0\end{pmatrix},
	\begin{pmatrix}1&0\\0&-1\end{pmatrix}
	\right\}
	\]

	se obtiene:

	\[
	\boxed{
	M_{B,B^*}(f)
	=
	\begin{pmatrix}
	1&0\\
	0&1\\
	0&0
	\end{pmatrix}
	}
	\]

	\[
	\boxed{
	M_{B^*,B''}(g)
	=
	\begin{pmatrix}
	1&0&1\\
	1&0&-1\\
	2&0&0
	\end{pmatrix}
	}
	\]

	y la matriz de $g\circ f$ permanece:

	\[
	\boxed{
	M_{B,B''}(g\circ f)
	=
	\begin{pmatrix}
	1&0\\
	1&0\\
	2&0
	\end{pmatrix}
	}
	\]

\section{Ejercicio 23}
	\subsection{Solución}

	Sea:

	\[
	p(x)=a+bx+cx^2
	\]

	Entonces:

	\[
	p'(x)=b+2cx
	\]

	\[
	p'(0)=b
	\]

	\[
	p'(1)=b+2c
	\]

	\[
	p'(-1)=b-2c
	\]

	Por tanto:

	\[
	f(a+bx+cx^2)
	=
	\begin{pmatrix}
	b&b+2c\\
	b+2c&b-2c
	\end{pmatrix}
	\]

	La expresión depende linealmente de $a,b,c$, luego:

	\[
	\boxed{
	f\text{ es lineal}
	}
	\]

	En las bases canónicas:

	\[
	\boxed{
	M(f)
	=
	\begin{pmatrix}
	0&1&0\\
	0&1&2\\
	0&1&-2
	\end{pmatrix}
	}
	\]

	El núcleo canónico es:

	\[
	\Ker f
	=
	\Span\{1\}
	\]

	En la base:

	\[
	B=\{x^2,(x-1)^2,(x+1)^2\}
	\]

	se tiene:

	\[
	1
	=
	-x^2
	+
	\frac12(x-1)^2
	+
	\frac12(x+1)^2
	\]

	Por tanto:

	\[
	[\Ker f]_B
	=
	\Span\left\{
	\left(-1,\frac12,\frac12\right)
	\right\}
	\]

	Equivalentemente:

	\[
	\boxed{
	\Ker f
	=
	\{(x,y,z)_B:y=z,\ x=-2z\}
	}
	\]

	Para:

	\[
	h=g\circ f
	\]

	se obtiene:

	\[
	h(a+bx+cx^2)
	=
	(2b-2c)+2cx
	\]

	Por tanto:

	\[
	M_{B_c}(h)
	=
	\begin{pmatrix}
	0&2&-2\\
	0&0&2
	\end{pmatrix}
	\]

	La condición:

	\[
	h(p)\in T=\Span\{1+x\}
	\]

	implica:

	\[
	2b-2c=2c
	\]

	\[
	b=2c
	\]

	Así:

	\[
	\boxed{
	h^{-1}(T)
	=
	\Span\{1,2x+x^2\}
	}
	\]

	En las bases:

	\[
	B=\{x^2,(x-1)^2,(x+1)^2\}
	\]

	y:

	\[
	B'=\{1-2x,x\}
	\]

	se obtiene:

	\[
	\boxed{
	M_{B,B'}(h)
	=
	\begin{pmatrix}
	-2&-6&2\\
	-2&-10&6
	\end{pmatrix}
	}
	\]

	Para:

	\[
	p(x)=3-2x+6x^2
	\]

	sus coordenadas en $B$ son:

	\[
	[p]_B=(3,2,1)
	\]

	Entonces:

	\[
	[h(p)]_{B'}
	=
	\begin{pmatrix}
	-2&-6&2\\
	-2&-10&6
	\end{pmatrix}
	\begin{pmatrix}
	3\\2\\1
	\end{pmatrix}
	=
	\begin{pmatrix}
	-16\\-20
	\end{pmatrix}
	\]

	Pasando a la base canónica:

	\[
	-16(1-2x)-20x
	=
	-16+12x
	\]

	\[
	\boxed{
	h(3-2x+6x^2)
	=
	-16+12x
	}
	\]

\end{document}
